\documentclass[10pt,a4paper]{amsart}
\usepackage{amsmath,amssymb,amsthm,graphicx}
% Retained preamble for the Stacks Project chapter snapshot in this directory.
% The chapter's own preamble.tex is a copy of the shared Stacks Project preamble; the block
% below is sliced MECHANICALLY, line for line, from it. Nothing was added, removed or
% redefined; the wrapper class, the xr-hyper/hyperref wiring and the externaldocument list
% are omitted because the assembler supplies the document class.
%
% lines taken from preamble.tex: 3, 5, 9, 11, 12, 156, 157, 158, ... (51 lines total)
%
% package coverage restored from the chapter's own preamble (2026-10-13): xy, xy-2cell

\usepackage[all]{xy}
\xyoption{2cell}
\UseAllTwocells

\newenvironment{reference}{\comment}{\endcomment}
\newenvironment{slogan}{\comment}{\endcomment}
\newenvironment{history}{\comment}{\endcomment}
\theoremstyle{plain}
\newtheorem{theorem}[subsection]{Theorem}
\newtheorem{proposition}[subsection]{Proposition}
\newtheorem{lemma}[subsection]{Lemma}
\theoremstyle{definition}
\newtheorem{definition}[subsection]{Definition}
\newtheorem{example}[subsection]{Example}
\newtheorem{exercise}[subsection]{Exercise}
\newtheorem{situation}[subsection]{Situation}
\theoremstyle{remark}
\newtheorem{remark}[subsection]{Remark}
\newtheorem{remarks}[subsection]{Remarks}
\numberwithin{equation}{subsection}
\def\lim{\mathop{\mathrm{lim}}\nolimits}
\def\colim{\mathop{\mathrm{colim}}\nolimits}
\def\Spec{\mathop{\mathrm{Spec}}}
\def\Hom{\mathop{\mathrm{Hom}}\nolimits}
\def\Ext{\mathop{\mathrm{Ext}}\nolimits}
\def\SheafHom{\mathop{\mathcal{H}\!\mathit{om}}\nolimits}
\def\SheafExt{\mathop{\mathcal{E}\!\mathit{xt}}\nolimits}
\def\Sch{\mathit{Sch}}
\def\Mor{\mathop{\mathrm{Mor}}\nolimits}
\def\Ob{\mathop{\mathrm{Ob}}\nolimits}
\def\Sh{\mathop{\mathit{Sh}}\nolimits}
\def\NL{\mathop{N\!L}\nolimits}
\def\CH{\mathop{\mathrm{CH}}\nolimits}
\def\proetale{{pro\text{-}\acute{e}tale}}
\def\etale{{\acute{e}tale}}
\def\QCoh{\mathit{QCoh}}
\def\Ker{\mathop{\mathrm{Ker}}}
\def\Im{\mathop{\mathrm{Im}}}
\def\Coker{\mathop{\mathrm{Coker}}}
\def\Coim{\mathop{\mathrm{Coim}}}
\DeclareMathSymbol{\boxtimes}{\mathbin}{AMSa}{"02}
\def\QCohstack{\mathcal{QC}\!\mathit{oh}}
\def\Cohstack{\mathcal{C}\!\mathit{oh}}
\def\Spacesstack{\mathcal{S}\!\mathit{paces}}
\def\Quotfunctor{\mathrm{Quot}}
\def\Hilbfunctor{\mathrm{Hilb}}
\def\Curvesstack{\mathcal{C}\!\mathit{urves}}
\def\Polarizedstack{\mathcal{P}\!\mathit{olarized}}
\def\Complexesstack{\mathcal{C}\!\mathit{omplexes}}
\def\Pic{\mathop{\mathrm{Pic}}\nolimits}
\def\Picardstack{\mathcal{P}\!\mathit{ic}}
\def\Picardfunctor{\mathrm{Pic}}
\def\Deformationcategory{\mathcal{D}\!\mathit{ef}}
\usepackage{cleveref}
\crefname{originallemma}{original lemma}{original lemmas}
\Crefname{originallemma}{Original lemma}{Original lemmas}
\makeatletter
\newlabel{algebra-definition-relative-global-complete-intersection}{{00SP}{source}{}{sourcepointer-1}{}}
\newlabel{algebra-definition-relative-global-complete-intersection@cref}{{[originallemma][0][]00SP}{[1][0][]source}{}{}{}}
\makeatother
\crefname{originallemma}{original lemma}{original lemmas}
\Crefname{originallemma}{Original lemma}{Original lemmas}
\makeatletter
\newlabel{algebra-lemma-relative-global-complete-intersection}{{00SW}{source}{}{sourcepointer-2}{}}
\newlabel{algebra-lemma-relative-global-complete-intersection@cref}{{[originallemma][0][]00SW}{[1][0][]source}{}{}{}}
\makeatother
\crefname{originallemma}{original lemma}{original lemmas}
\Crefname{originallemma}{Original lemma}{Original lemmas}
\makeatletter
\newlabel{lemma-tate}{{0BWC}{source}{}{sourcepointer-3}{}}
\newlabel{lemma-tate@cref}{{[originallemma][0][]0BWC}{[1][0][]source}{}{}{}}
\makeatother
\crefname{originallemma}{original lemma}{original lemmas}
\Crefname{originallemma}{Original lemma}{Original lemmas}
\makeatletter
\newlabel{more-algebra-lemma-koszul-regular-flat-base-change}{{062H}{source}{}{sourcepointer-4}{}}
\newlabel{more-algebra-lemma-koszul-regular-flat-base-change@cref}{{[originallemma][0][]062H}{[1][0][]source}{}{}{}}
\makeatother
\crefname{originallemma}{original lemma}{original lemmas}
\Crefname{originallemma}{Original lemma}{Original lemmas}
\makeatletter
\newlabel{more-algebra-lemma-relative-global-complete-intersection-koszul}{{07D2}{source}{}{sourcepointer-5}{}}
\newlabel{more-algebra-lemma-relative-global-complete-intersection-koszul@cref}{{[originallemma][0][]07D2}{[1][0][]source}{}{}{}}
\makeatother
\crefname{originallemma}{original lemma}{original lemmas}
\Crefname{originallemma}{Original lemma}{Original lemmas}
\makeatletter
\newlabel{more-algebra-lemma-truncate-koszul-regular}{{068M}{source}{}{sourcepointer-6}{}}
\newlabel{more-algebra-lemma-truncate-koszul-regular@cref}{{[originallemma][0][]068M}{[1][0][]source}{}{}{}}
\makeatother
\title{Stacks, Discriminants: lemma-quasi-finite-complete-intersection}
\author{The Stacks Project Authors}
\date{stacks-project:discriminant.tex:local-snapshot-20261012; GNU FDL 1.2 or later}
\begin{document}
\maketitle
\noindent\textit{Editorial note.} Statement span is lines 2611-2626 of sources/stacks-discriminant/pinned-original-source.tex; proof span is lines 2628-2705, the complete author proof (76 body lines).\par\medskip
\noindent\textit{Editorial note.} Navigation stubs: 6 cited-but-not-included labels are supplied as non-load-bearing pointers with their pinned tags where the snapshot tag index has them.\par\medskip
\noindent\textit{Editorial note.} extra\_packages is empty on purpose: the retained preamble already loads xy with options and a second declaration clashes in the assembler wrapper.\par\medskip
\noindent\textit{Editorial note.} Text is retained verbatim, including the author's own typos and placeholders.\par\medskip
\noindent\textit{Editorial note.} Queue index 2 of queue\_discriminant.json (pb=3185, stmt 2611-2626, proof 2628-2705).\par\medskip
\noindent\textit{Allowed standard prerequisites (editorial).}\begin{itemize}
\item Explicit prerequisite boundary. The cited inputs are prerequisites; the argument in the retained span is not replaced by any cited result.
\end{itemize}\medskip
\ifdefined\hypertarget\hypertarget{sourcepointer-1}{}\fi
\noindent\textit{Source reference.} Original lemma 00SP: sources/stacks-discriminant/pinned-original-source.tex, source lines 2705--2705. Cited by the retained statement or proof but not part of the collected spans; supplied only as a navigation pointer. The argument does not depend on its content.\par\medskip
\ifdefined\hypertarget\hypertarget{sourcepointer-2}{}\fi
\noindent\textit{Source reference.} Original lemma 00SW: sources/stacks-discriminant/pinned-original-source.tex, source lines 2705--2705. Cited by the retained statement or proof but not part of the collected spans; supplied only as a navigation pointer. The argument does not depend on its content.\par\medskip
\ifdefined\hypertarget\hypertarget{sourcepointer-3}{}\fi
\noindent\textit{Source reference.} Original lemma 0BWC: sources/stacks-discriminant/pinned-original-source.tex, source lines 2573--2576. Cited by the retained statement or proof but not part of the collected spans; supplied only as a navigation pointer. The argument does not depend on its content.\par\medskip
\ifdefined\hypertarget\hypertarget{sourcepointer-4}{}\fi
\noindent\textit{Source reference.} Original lemma 062H: sources/stacks-discriminant/pinned-original-source.tex, source lines 2705--2705. Cited by the retained statement or proof but not part of the collected spans; supplied only as a navigation pointer. The argument does not depend on its content.\par\medskip
\ifdefined\hypertarget\hypertarget{sourcepointer-5}{}\fi
\noindent\textit{Source reference.} Original lemma 07D2: sources/stacks-discriminant/pinned-original-source.tex, source lines 2705--2705. Cited by the retained statement or proof but not part of the collected spans; supplied only as a navigation pointer. The argument does not depend on its content.\par\medskip
\ifdefined\hypertarget\hypertarget{sourcepointer-6}{}\fi
\noindent\textit{Source reference.} Original lemma 068M: sources/stacks-discriminant/pinned-original-source.tex, source lines 2705--2705. Cited by the retained statement or proof but not part of the collected spans; supplied only as a navigation pointer. The argument does not depend on its content.\par\medskip
\begin{lemma}
\label{lemma-quasi-finite-complete-intersection}
Let $A$ be a ring. Let $n \geq 1$ and
$h, f_1, \ldots, f_n \in A[x_1, \ldots, x_n]$.
Set $B = A[x_1, \ldots, x_n, 1/h]/(f_1, \ldots, f_n)$.
Assume that $B$ is quasi-finite over $A$.
Then
\begin{enumerate}
\item $B$ is flat over $A$ and $A \to B$ is a relative local complete
intersection,
\item the annihilator $J$ of $I = \Ker(B \otimes_A B \to B)$
is free of rank $1$ over $B$,
\item the Noether different of $B$ over $A$ is generated
by $\det(\partial f_i/\partial x_j)$ in $B$.
\end{enumerate}
\end{lemma}

\begin{proof}
Note that
$B = A[x, x_1, \ldots, x_n]/(xh - 1, f_1, \ldots, f_n)$
is a relative global complete intersection over $A$, see
Algebra, Definition
\ref{algebra-definition-relative-global-complete-intersection}.
By Algebra, Lemma \ref{algebra-lemma-relative-global-complete-intersection}
we see that $B$ is flat over $A$.

\medskip\noindent
Write $P' = A[x, x_1, \ldots, x_n]$ and
$P = P'/(xh - 1) = A[x_1, \ldots, x_n, 1/h]$.
Then we have $P' \to P \to B$.
By More on Algebra, Lemma
\ref{more-algebra-lemma-relative-global-complete-intersection-koszul}
we see that $xh - 1, f_1, \ldots, f_n$ is a Koszul regular sequence
in $P'$. Since $xh - 1$ is a Koszul regular sequence of length
one in $P'$ (by the same lemma for example) we conclude that 
$f_1, \ldots, f_n$ is a Koszul regular sequence in $P$ by
More on Algebra, Lemma \ref{more-algebra-lemma-truncate-koszul-regular}.

\medskip\noindent
Let $g_i \in P \otimes_A B$ be the image of $x_i \otimes 1 - 1 \otimes x_i$.
Let us use the short hand $y_i = x_i \otimes 1$ and $z_i = 1 \otimes x_i$
in $A[x_1, \ldots, x_n] \otimes_A A[x_1, \ldots, x_n]$
so that $g_i$ is the image of $y_i - z_i$. For a polynomial
$f \in A[x_1, \ldots, x_n]$ we write $f(y) = f \otimes 1$
and $f(z) = 1 \otimes f$ in the above tensor product.
Then we have
$$
P \otimes_A B/(g_1, \ldots, g_n) =
\frac{A[y_1, \ldots, y_n, z_1, \ldots, z_n, \frac{1}{h(y)h(z)}]}
{(f_1(z), \ldots, f_n(z), y_1 - z_1, \ldots, y_n - z_n)}
$$
which is clearly isomorphic to $B$. Hence by the same arguments
as above we find that $f_1(z), \ldots, f_n(z), y_1 - z_1, \ldots, y_n - z_n$
is a Koszul regular sequence in
$A[y_1, \ldots, y_n, z_1, \ldots, z_n, \frac{1}{h(y)h(z)}]$.
The sequence $f_1(z), \ldots, f_n(z)$ is a Koszul regular in
$A[y_1, \ldots, y_n, z_1, \ldots, z_n, \frac{1}{h(y)h(z)}]$
by flatness of the map
$$
P \longrightarrow A[y_1, \ldots, y_n, z_1, \ldots, z_n,
\textstyle{\frac{1}{h(y)h(z)}}],\quad x_i \longmapsto z_i
$$
and More on Algebra, Lemma
\ref{more-algebra-lemma-koszul-regular-flat-base-change}.
By More on Algebra, Lemma \ref{more-algebra-lemma-truncate-koszul-regular}
we conclude that $g_1, \ldots, g_n$ is a regular sequence
in $P \otimes_A B$.

\medskip\noindent
At this point we have verified all the assumptions of Lemma \ref{lemma-tate}
above with $P$, $f_1, \ldots, f_n$, and $g_i \in P \otimes_A B$ as above.
In particular the annihilator $J$ of $I$ is freely generated by one
element $\delta$ over $B$.
Set $f_{ij} = \partial f_i/\partial x_j \in A[x_1, \ldots, x_n]$.
An elementary computation shows that we can write
$$
f_i(y) =
f_i(z_1 + g_1, \ldots, z_n + g_n) =
f_i(z) + \sum\nolimits_j f_{ij}(z) g_j +
\sum\nolimits_{j, j'} F_{ijj'}g_jg_{j'}
$$
for some $F_{ijj'} \in A[y_1, \ldots, y_n, z_1, \ldots, z_n]$.
Taking the image in $P \otimes_A B$ the terms $f_i(z)$ map to
zero and we obtain
$$
f_i \otimes 1 = \sum\nolimits_j
\left(1 \otimes f_{ij} + \sum\nolimits_{j'} F_{ijj'}g_{j'}\right)g_j
$$
Thus we conclude from Lemma \ref{lemma-tate}
that $\delta = \det(g_{ij})$ with
$g_{ij} = 1 \otimes f_{ij} + \sum_{j'} F_{ijj'}g_{j'}$.
Since $g_{j'}$ maps to zero in $B$, we conclude
that the image of $\det(\partial f_i/\partial x_j)$ in $B$
generates the Noether different of $B$ over $A$.
\end{proof}

\end{document}
